\chapter{Matrices}

These definitions are just for a quick review. Please refer to Artin's Algebra Chapter 1.1-1.3 for more detailed information.

\section{Definitions}

\begin{definitionenv}
	Let $\mathcal{K}$ be a field and let $m, n$ be two positive integers.
	An array of elements in $\mathcal{K}$ of the form
	\[
	A =
	\begin{pmatrix}
		a_{11} & a_{12} & \cdots & a_{1n} \\
		a_{21} & a_{22} & \cdots & a_{2n} \\
		\vdots & \vdots & \ddots & \vdots \\
		a_{m1} & a_{m2} & \cdots & a_{mn}
	\end{pmatrix}
	\]
	is called a \textbf{matrix} over $\mathcal{K}$ of size $m \times n$, or an $m$-by-$n$ matrix.
\end{definitionenv}

\begin{definitionenv}
	For convenience, we may write $A = (a_{ij})$ for $1 \leq i \leq m$ and $1 \leq j \leq n$ to denote a matrix in $\mathcal{K}^{m \times n}$.
	The element $a_{ij}$ is called the $(i,j)$-entry of $A$.
	Each \textbf{column} of $A$ is a vector in $\mathcal{K}^m$, and each \textbf{row} of $A$ is a vector in $\mathcal{K}^n$.
\end{definitionenv}


\begin{definitionenv}
	If $m = n$, the matrix $A$ is called a \textbf{square matrix}.
	If $m \neq n$, $A$ is called a \textbf{non-square matrix}.
	The \textbf{zero matrix} is the matrix all of whose entries are zero.
\end{definitionenv}

\begin{definitionenv}
	Let $A = (a_{ij})$ and $B = (b_{ij})$ be two matrices in $\M_{m \times n}(\mathcal{K})$, and let $\alpha \in \mathcal{K}$.
	Matrix addition and scalar multiplication are defined as:
	\[
	(A + B) = (a_{ij} + b_{ij}), \quad
	(\alpha A) = (\alpha a_{ij}), \quad
	\text{for } 1 \leq i \leq m, \ 1 \leq j \leq n.
	\]
\end{definitionenv}

\begin{remarkenv}
	Let $\M_{m \times n}(\mathcal{K})$ (or $\mathrm{Mat}_{m \times n}(\mathcal{K})$) denote the set of all $m \times n$ matrices over $\mathcal{K}$.
	Equipped with the matrix addition and scalar multiplication defined above,
	$\M_{m \times n}(\mathcal{K})$ forms a \textbf{vector space} over $\mathcal{K}$.

	Its dimension is
	\[
	\dim \M_{m \times n}(\mathcal{K}) = m n.
	\]
\end{remarkenv}

\begin{definitionenv}
	The \textbf{standard basis} of the vector space $\M_{m \times n}(\mathcal{K})$
	is the collection of $m n$ matrices $\{e_{ij}\}$, where each $e_{ij}$ has entry
	\[
	(e_{ij})_{pq} =
	\begin{cases}
		1, & \text{if } p = i \text{ and } q = j,\\
		0, & \text{otherwise}.
	\end{cases}
	\]
	Each $e_{ij}$ has a single $1$ in the $(i,j)$-position and zeros elsewhere.
	Every matrix $A = (a_{ij}) \in \M_{m \times n}(\mathcal{K})$ can be expressed uniquely as
	\[
	A = \sum_{i=1}^{m} \sum_{j=1}^{n} a_{ij} e_{ij}.
	\]
\end{definitionenv}

\begin{definitionenv}
	Let $A = (a_{ij}) \in \M_{m \times n}(\mathcal{K})$ and
	$B = (b_{jk}) \in \M_{n \times p}(\mathcal{K})$.
	The \textbf{product} of $A$ and $B$, denoted by $AB$, is defined as the $m \times p$ matrix
	\[
	AB = (c_{ik}) \in \M_{m \times p}(\mathcal{K}),
	\]
	where each entry $c_{ik}$ is given by
	\[
	c_{ik} = \sum_{j=1}^{n} a_{ij} b_{jk},
	\quad \text{for } 1 \leq i \leq m, \ 1 \leq k \leq p.
	\]
\end{definitionenv}

\begin{theoremenv}
	(\textbf{Associative Law})
	Let $A \in \M_{m \times n}(\mathcal{K})$,
	$B \in \M_{n \times p}(\mathcal{K})$,
	and $C \in \M_{p \times q}(\mathcal{K})$.
	Then matrix multiplication is associative:
	\[
	(AB)C = A(BC).
	\]
\end{theoremenv}

\begin{theoremenv}
	(\textbf{Distributive Laws})
	For matrices $A, A_1, A_2 \in \M_{m \times n}(\mathcal{K})$
	and $B, B_1, B_2 \in \M_{n \times p}(\mathcal{K})$,
	matrix multiplication satisfies
	\[
	A(B_1 + B_2) = AB_1 + AB_2,
	\quad (A_1 + A_2)B = A_1B + A_2B.
	\]
\end{theoremenv}

\begin{definitionenv}
	For each positive integer $n$, the \textbf{identity matrix} of order $n$, denoted by $I_n$,
	is the $n \times n$ matrix whose entries are defined by
	\[
	(I_n)_{ij} =
	\begin{cases}
		1, & \text{if } i = j,\\
		0, & \text{otherwise}.
	\end{cases}
	\]
	The matrix $I_n$ satisfies the property that for any
	$A \in \M_{m \times n}(\mathcal{K})$ and
	$B \in \M_{n \times p}(\mathcal{K})$,
	\[
	AI_n = A, \quad I_nB = B.
	\]
\end{definitionenv}

\begin{definitionenv}
	Let $A \in \M_{n \times n}(\mathcal{K})$ be a square matrix.
	If there exists a matrix $B \in \M_{n \times n}(\mathcal{K})$ such that
	\[
	AB = BA = I_n,
	\]
	then $A$ is called an \textbf{invertible matrix} (or \textbf{non-singular matrix}),
	and $B$ is called the \textbf{inverse matrix} of $A$, denoted by $A^{-1}$. If no such $B$ exists, $A$ is called a \textbf{singular matrix}.
\end{definitionenv}

\begin{remarkenv}
	Such inverse shown above is unique.
\end{remarkenv}

\begin{definitionenv}
	Let $A = (a_{ij}) \in \M_{m \times n}(\mathcal{K})$.
	The \textbf{transpose} of $A$, denoted by $A^{T}$, is the $n \times m$ matrix obtained by interchanging the rows and columns of $A$.
	Formally,
	\[
	A^{T} = (a_{ji}) \in \M_{n \times m}(\mathcal{K}),
	\]
	that is, the $(i,j)$-entry of $A^{T}$ equals the $(j,i)$-entry of $A$. If $A^{T} = A$, then $A$ is called a \textbf{symmetric matrix}.
\end{definitionenv}

\begin{theoremenv}
	(\textbf{Properties of the Transpose})
	Let $A, B \in \M_{m \times n}(\mathcal{K})$ and $\alpha \in \mathcal{K}$.
	Then the transpose operation satisfies the following properties:
	\begin{enumerate}
		\item $(A^{T})^{T} = A$;
		\item $(A + B)^{T} = A^{T} + B^{T}$;
		\item $(\alpha A)^{T} = \alpha A^{T}$;
		\item If $A \in \M_{m \times n}(\mathcal{K})$ and
		$B \in \M_{n \times p}(\mathcal{K})$, then
		\[
		(AB)^{T} = B^{T}A^{T}.
		\]
	\end{enumerate}
\end{theoremenv}


\section{Linear Equations}
\begin{definitionenv}
	A \textbf{linear system} over a field $\mathcal{K}$ is a collection of $m$ linear equations in $n$ variables or unknowns
	$x_1, x_2, \dots, x_n$ of the form
	\[
	\begin{cases}
		a_{11}x_1 + a_{12}x_2 + \cdots + a_{1n}x_n = b_1,\\
		a_{21}x_1 + a_{22}x_2 + \cdots + a_{2n}x_n = b_2,\\
		\quad \vdots\\
		a_{m1}x_1 + a_{m2}x_2 + \cdots + a_{mn}x_n = b_m,
	\end{cases}
	\]
	where $a_{ij}, b_i \in \mathcal{K}$ for all $i,j$.

	It can be written compactly in matrix form as
	\[
	A\bm{x} = \bm{b},
	\]
	where
	\[
	A =
	\begin{pmatrix}
		a_{11} & a_{12} & \cdots & a_{1n} \\
		a_{21} & a_{22} & \cdots & a_{2n} \\
		\vdots & \vdots & \ddots & \vdots \\
		a_{m1} & a_{m2} & \cdots & a_{mn}
	\end{pmatrix},
	\quad
	\bm{x} =
	\begin{pmatrix}
		x_1 \\ x_2 \\ \vdots \\ x_n
	\end{pmatrix},
	\quad
	\bm{b} =
	\begin{pmatrix}
		b_1 \\ b_2 \\ \vdots \\ b_m
	\end{pmatrix}.
	\]
\end{definitionenv}


\begin{definitionenv}
	A linear system $A\bm{x} = \bm{b}$ is called:
	\begin{itemize}
		\item \textbf{homogeneous} if $\bm{b} = \bm{0}$, i.e. the system is of the form $A\bm{x} = \bm{0}$;
		\item \textbf{inhomogeneous} (or nonhomogeneous) if $\bm{b} \neq \bm{0}$.
	\end{itemize}
\end{definitionenv}

\begin{definitionenv}
	A vector $\bm{x} \in \mathcal{K}^n$ satisfying the equation $A\bm{x} = \bm{b}$ is called a \textbf{solution} of the system.

	In the special case of a homogeneous system $A\bm{x} = \bm{0}$:
	\begin{itemize}
		\item the vector $\bm{0}$ is called the \textbf{trivial solution};
		\item any nonzero vector $\bm{x} \neq \bm{0}$ satisfying $A\bm{x} = \bm{0}$ is called a \textbf{nontrivial solution}.
	\end{itemize}
\end{definitionenv}

\begin{theoremenv}
	Let
	\[
	\begin{cases}
		a_{11}x_1 + a_{12}x_2 + \cdots + a_{1n}x_n = 0,\\
		a_{21}x_1 + a_{22}x_2 + \cdots + a_{2n}x_n = 0,\\
		\quad \vdots \\
		a_{m1}x_1 + a_{m2}x_2 + \cdots + a_{mn}x_n = 0,
	\end{cases}
	\]
	be a homogeneous system of $m$ linear equations in $n$ unknowns, with coefficients in a field $\mathcal{K}$.
	Assume that $n > m$.
	Then the system has a \textbf{nontrivial solution} in $\mathcal{K}$.
\end{theoremenv}

\begin{theoremenv}
	Assume that $m = n$ in the system above, and that the column vectors
	$A^1, A^2, \dots, A^n$ of the coefficient matrix $A$ are linearly independent.
	Then the system has a \textbf{unique solution} in $\mathcal{K}$.
\end{theoremenv}

\newpage